\paragraph{}
Tato sekce shrnuje pět klíčových etických principů, které by měly řídit nasazení generativní AI ve výuce: transparentnost, ochrana soukromí, spravedlnost (equita), akademická integrita a zodpovědnost. Následující stručné definice a praktické poznámky slouží jako pracovní rámec pro garanty předmětů, vyučující a metodiky. Každý princip obsahuje krátké vysvětlení, konkrétní doporučení pro praxi a upozornění na běžné rizikové situace, které se v pedagogickém prostředí mohou vyskytnout.

\subparagraph{Transparentnost}
- Definice (general background knowledge): jasné informování studentů o tom, kde a jaký nástroj AI je používán, jaké jsou limity jeho výstupů a jak mají studenti deklarovat své využití AI. Transparentnost zahrnuje i sdílení informací o tom, zda byly použity předpřipravené šablony, externí dataset nebo modely třetích stran.
- Praktická doporučení: vložte krátké oznámení o použití AI do sylabu, uveďte verzi/nastavení asistenta a pravidla pro citování AI‑výstupů (viz příslušné šablony v přílohách). Uveďte, jak budou výsledky AI ověřovány a jaký vliv může mít použití AI na hodnocení. Doporučujeme také informovat studenty o možných chybách modelu a poskytnout krátký návod, jak ověřit faktické tvrzení v generovaném textu. (general background knowledge)

\subparagraph{Ochrana soukromí}
- Definice (general background knowledge): minimalizace sběru osobních údajů, pseudonymizace tam, kde je to možné, a předchozí posouzení dopadů (DPIA) pro nové piloty. Zahrnuje také posouzení přenosu dat mezi institucí a poskytovateli cloudových služeb a dodržování relevantních právních předpisů.
- Praktická implikace: při používání edukačních nástrojů, které sbírají postupy a výsledky žáků (např. nástroje typu „Pokrok v matematice“), ověřte, jaká data jsou odesílána a kde jsou uložena \cite{kb_ai_101_tipu_pdf_04e4a115_page_8_1}. Doporučujeme omezit nahrávání citlivých nebo přímo identifikovatelných údajů do externích služeb. Před spuštěním pilotu dokumentujte, jaká metadata se ukládají, kdo má k nim přístup, a stanovte dobu uchovávání dat. (general background knowledge)

\subparagraph{Spravedlnost / equita}
- Definice (general background knowledge): návrh aktivit a podpory tak, aby přístup k AI ani jeho výsledky nezvýhodňovaly systematicky některé skupiny studentů. Princip se vztahuje jak na sémantickou stránku výstupů (obsah a jazyk), tak na technický přístup (přístup k hardware a síťovým zdrojům).
- Praktické kroky: zajistěte alternativní způsoby přístupu (offline materiály, asistivní technologie) a pravidelně kontrolujte, zda generované materiály nepřevádějí nebo nezesilují existující biasy. Při demonstracích využijte ověřené edukační prostředí (např. Minecraft Education pro názorné ukázky principů) s ohledem na dostupnost a licencování \cite{kb_ai_101_tipu_pdf_04e4a115_page_15_2}. Provádějte testování materiálů na reprezentativním vzorku studentů a zaznamenávejte zjištěné disproporce pro následné úpravy obsahů.

\subparagraph{Akademická integrita}
- Definice (general background knowledge): jasné normy toho, co je povoleno, povinnosti deklarace použití AI a hodnocení procesu místo pouhého produktu. Dále sem patří pravidla pro spoluvlastnictví výsledků a uznání přínosu nástrojů AI v tvůrčích úkolech.
- Praktická doporučení: přepracujte zadání tak, aby hodnotila proces (drafty, záznamy postupu, interview), a zvažte integraci polí pro deklaraci AI do odevzdávacích šablon. Pokud povolujete AI nástroje, požadujte reflexi o tom, jak byl nástroj použit. Doporučuje se také používat techniky ověřování originality a kontrolních otázek, které odhalí, zda student rozumí obsahu, který prezentuje. (general background knowledge)

\subparagraph{Zodpovědnost}
- Definice (general background knowledge): jasné určení, kdo nese odpovědnost za validaci výstupů AI (vyučující / garant / IT) a jak budou řešeny chybné či škodlivé výstupy. Zahrnuje také procedury pro hlášení incidentů a postupy nápravy v případě škodlivého či nevhodného obsahu.
- Doporučení pro provoz: vždy fungujte s lidským validátorem (human‑in‑the‑loop) pro kritické rozhodnutí; dokumentujte prompty a verze modelů pro případ auditů. Pokud používáte enterprise nástroje pro přípravu testů či materiálů (např. Microsoft 365 Copilot pro generování testů ve Forms), definujte roli lidského přezkumu při konečné editaci \cite{kb_ai_101_tipu_pdf_04e4a115_page_36_1}. Vypracujte jasný eskalační scénář pro situace, kdy výstup AI obsahuje chybná data nebo eticky problematický obsah, a určete kontaktní osoby zodpovědné za rychlou nápravu.

\paragraph{Krátký praktický checklist před pilotem (shrnutí, general background knowledge)}
\begin{itemize}
  \item Vymezte pedagogický cíl pilotu a očekávaný přínos.
  \item Proveďte základní posouzení rizik a DPIA (pokud jsou zpracovávána osobní data).
  \item Ověřte licenční a datové podmínky nástroje (kde jsou data uložena, kdo je zpracovává).
  \item Zajistěte human‑in‑the‑loop pro finální validaci citlivých výstupů.
  \item Připravte krátkou formulaci pro sylabus/informaci pro studenty o pravidlech použití AI.
  \item Zdokumentujte použité prompty, verze modelů a testovací scénáře pro budoucí revizi.
  \item Naplánujte monitorování pilotu (sběr zpětné vazby od studentů, metriky úspěšnosti).
  \item Určete odpovědné osoby a kontaktní postup pro případ bezpečnostních nebo etických incidentů.
\end{itemize}

\paragraph{}
Pro rychlé praktické příklady a šablony (označené texty do sylabů, checklisty DPIA a prompt‑šablony) využijte centrální repozitář VUT (ai.vut.cz), kde budou dostupné i video‑tutoriály a aktualizované materiály pro pedagogy \cite{kb_ai_101_tipu_pdf_04e4a115_page_15_2}. Repozitář obsahuje rovněž doporučené formulace pro informování studentů, vzory pro záznam použitých promptů a příklady jednoduchých DPIA formulářů, které lze upravit pro konkrétní piloty. Doporučujeme pravidelně kontrolovat aktualizace repozitáře před každým větším nasazením.